% Reference entries, P (AN-11).  Law 1: os/iword/core/edit.c iw_ed_hpage
% (Ctrl-Enter: the hard page code 0x0C, WP's D4/00 group kept before it by
% iw_ed_fixup; deleted without a question), face_a.c line_rows / paint_doc
% (the bar of = on a row of its own below the line), iw_ed_where (54 lines
% to a page; Pg), the File menu (face_a.c file_items: Retrieve, Save, Exit
% with ^Q; the TBX_EVT_MENU handler).
% GATES: test-pubs-iword test-iword-codes test-iword-wp5-roundtrip test-iword-face test-flair-iword-block
\Gates{test-pubs-iword test-iword-codes test-iword-wp5-roundtrip test-iword-face test-flair-iword-block}
\begin{Feature}{Page Break, Hard}
A hard page ends a page where you choose: the text after it starts on a new
page whatever room is left on this one.  Use it to start a new section, or
to keep a table or a signature block on a page of its own.

To insert a hard page,
\begin{WPSteps}
\item Move the cursor to where the new page is to begin.
\item Press \K{Ctrl-Enter}.
\end{WPSteps}
\Documents{wpkey}{Ctrl-Enter}

\BeginReplay{ref-hpage}
On the editing screen a hard page shows as a line of equal signs across
the screen below the line it ends.  For example, type
\TypeRun{Page one ends here.}, press \PressKey{Ctrl-Enter} and type
\TypeRun{Page two.}  The status line shows \Scr{Pg 2}:
\WordScreen[rows=3]{ref-hpage}{A hard page on the editing screen}
\EndReplay

\Notes
\Topic{Code}
A hard page is the code \Code{[HPg]}.  To remove a hard page, delete the
code: put the cursor at the start of the line after the bar and press
Backspace, or use Reveal Codes.  InitechWord deletes it without asking.

\Topic{Pages Without Hard Pages}
A page holds 54 lines of single-spaced text.  After 54 lines the status line
counts the next line as the first line of the next page, \Scr{Ln 1"}; the
editing screen shows no line between the pages.  A hard page starts a new
page before that.

\Topic{The Bar}
The bar of equal signs is not part of the text.  The cursor never rests on
it, and it is not saved.

\SeeAlso{Move, Sentence/Paragraph/Page; Return, Hard; Status Line}
\end{Feature}

\begin{Feature}{Page Break, Soft}
InitechWord counts the lines of a document into pages of 54 single-spaced
lines, the lines that fit between the one inch top and bottom margins of an
11 inch page.  The status line shows the page and line the cursor is on;
there is no code for the end of such a page, and the editing screen draws no
line between the pages.

\BeginReplay{ref-softpage}
\Typing{Line 1}\Keys{Enter}%
For example, in a document of 56 short lines, the 55th line is the first
line of page 2:
\directlua{for i = 2, 55 do tex.sprint("\string\\Typing{Line " .. i .. "}\string\\Keys{Enter}") end}%
\Typing{Line 56}\Keys{Up}%
\WordScreen[rows=23]{ref-softpage}{Line 55 is on page 2, at Ln 1"}
\EndReplay

\Notes
\Topic{Hard Pages}
To end a page sooner, put in a hard page (see \InRef{Page Break, Hard}).
A hard page starts the count of the next page again.

\Topic{Move, Page}
Move's Page selection uses these pages (see \InRef{Move,
Sentence/Paragraph/Page}).

\SeeAlso{Move, Sentence/Paragraph/Page; Page Break, Hard; Status Line}
\end{Feature}

\begin{Feature}{Pull-Down Menus}
InitechWord's pull-down menu is the \textbf{File} menu, in the second row
of the menu bar while InitechWord is the application in front.  It holds
three items, each of which starts the same feature as a function key.

To choose an item from the File menu,
\begin{WPSteps}
\item Point to \textbf{File} in the menu bar and hold down the mouse button.
The menu drops down.
\item Drag down to the item you want and release the button.
\end{WPSteps}
\Documents{wppull}{File,Retrieve}
\Documents{wppull}{File,Save}
\Documents{wppull}{File,Exit}
\par\addvspace{4pt}
\noindent\begin{tabular}{@{}lll@{}}
\textbf{Item} & \textbf{Same as} & \textbf{See}\\
Retrieve & \FK{Retrieve}{Shift-F10} & \InRef{Retrieve}\\
Save & \FK{Save}{F10} & \InRef{Save}\\
Exit & \FK{Exit}{F7} & \InRef{Exit}\\
\end{tabular}
\par\addvspace{4pt}

\Notes
\Topic{Exit and Its Keyboard Equivalent}
The menu shows Exit with the command key and Q.  On your keyboard the
command key is Ctrl: hold down Ctrl and press Q.  Exit from the menu ends
a document you have not changed since you last saved or retrieved it
without asking; otherwise it asks the questions of Exit (F7).

\Topic{When the Menu Is Not Available}
The File menu items act only while the editing screen is shown.  While a
prompt or a menu of InitechWord is on the status line, or a Format screen is
shown, choosing them does nothing; answer or cancel the prompt first.

\SeeAlso{Exit; Retrieve; Save}
\end{Feature}
